\subsection{实验设置}
\label{sec:experiment-setup}

激光超声检测系统由两个主要部分组成：超声波激发子系统和超声波检测子系统。在超声波激发方面，采用调Q Nd:YAG激光器工作于基频波长 ($\lambda$ = 1064\,nm)。激光脉冲持续时间为纳秒量级，通过试件表面的光热效应产生宽带超声波。在检测方面，采用激光多普勒振动仪(LDV)测量试件表面的面外速度。在部分配置中，带宽合适的压电传感器(PZT)作为替代检测方法使用。

扫描配置采用计算机控制的平移台，使激发激光或试件相对于检测探头移动。评估了两种扫描模式：逐点机械扫描用于离散测量位置，以及沿试件表面预定路径的线扫描。

本研究使用了两块铝板试件。训练试件为完整铝板，标称尺寸为300\,mm $\times$ 300\,mm $\times$ 3\,mm（长$\times$宽$\times$厚）。测试试件为具有相同标称尺寸的铝板，包含一系列不同深度和宽度的机加工槽和缺口，作为受控缺陷目标用于评估检测能力。

数据采集使用数字示波器或等效高速采集系统。采样率设置为至少为有效带宽的两倍以满足奈奎斯特准则。采集带宽配置为捕获超声波信号的完整频率内容。采用信号平均法提高信噪比，平均次数根据所需的噪声降低和采集时间约束进行调整。

具体参数——包括激光脉冲能量、试件表面的光斑尺寸、扫描步长和精确检测带宽——以占位符形式提供于系统配置中，将在实验执行前确定。

% 占位符：激光脉冲能量（例如 50--200\,mJ）
% 占位符：激光光斑直径（例如 3--6\,mm）
% 占位符：扫描步长（例如 1--5\,mm）
% 占位符：LDV检测带宽（例如 DC至20\,MHz）
% 占位符：平均次数（例如 32、64、128）